% Generated draft: requires mathematical and source review.
\documentclass[11pt]{article}
\usepackage[margin=25mm]{geometry}
\usepackage{fontspec}
\setmainfont{Latin Modern Roman}
\usepackage{amsmath,amssymb,mathtools,unicode-math}
\setmathfont{Latin Modern Math}
\usepackage{enumitem,needspace,xurl,hyperref}
\setlength{\parindent}{0pt}
\setlength{\parskip}{5pt}
\setlength{\emergencystretch}{3em}
\newcommand{\N}{\mathbb{N}}
\newcommand{\Z}{\mathbb{Z}}
\newcommand{\Q}{\mathbb{Q}}
\newcommand{\R}{\mathbb{R}}
\newcommand{\C}{\mathbb{C}}
\newcommand{\F}{\mathbb{F}}
\newcommand{\A}{\mathbb{A}}
\newcommand{\PP}{\mathbb{P}}
\newcommand{\E}{\mathbb{E}}
\newcommand{\eps}{\varepsilon}
\newcommand{\dd}{\,\mathrm{d}}
\begin{document}
\section*{UnsolvedMath ID: 30002518}
\textbf{Source problem code:} catalogue entry
\subsection{Definitions and mathematical statement}
Let \(\mathcal A=\mathcal A(\mathbf x,B)\) denote a cluster algebra with initial cluster \(\mathbf x\) and exchange data \(B\). A cluster algebra is generated over its coefficient ring by the cluster variables obtained through seed mutation. The supplied record does not specify the coefficient ring or coefficient system, the permitted exchange data, or the convention for frozen variables and their invertibility. These choices remain unresolved in the quantifier ``every cluster algebra.''

An abelian group \(G\), written additively, is torsion-free if
\[
\forall g\in G\;\forall m\in\Z_{>0},\qquad mg=0\Longrightarrow g=0.
\]
For a Krull domain \(\mathcal A\) with fraction field \(K\), its divisor class group has the standard description
\[
\operatorname{Div}(\mathcal A)
 =\bigoplus_{\substack{\mathfrak p\in\operatorname{Spec}(\mathcal A)\\\operatorname{ht}(\mathfrak p)=1}}\Z[\mathfrak p],
\qquad
\operatorname{Cl}(\mathcal A)
 =\operatorname{Div}(\mathcal A)\big/\operatorname{Prin}(\mathcal A),
\]
where
\[
\operatorname{Prin}(\mathcal A)
 =\{\operatorname{div}(f):f\in K^{\times}\},
\qquad
\operatorname{div}(f)
 =\sum_{\operatorname{ht}(\mathfrak p)=1}v_{\mathfrak p}(f)[\mathfrak p].
\]
Here \(v_{\mathfrak p}\) is the normalized discrete valuation associated with the discrete valuation ring \(\mathcal A_{\mathfrak p}\); the Krull-domain hypothesis ensures finite support of each principal divisor.

\textbf{Conjecture as recorded.} The divisor class group of every cluster algebra is torsion-free. Explicitly, the intended assertion is
\[
\forall\mathcal A=\mathcal A(\mathbf x,B)\;\forall c\in\operatorname{Cl}(\mathcal A)\;\forall m\in\Z_{>0},
\qquad mc=0\Longrightarrow c=0,
\]
where the first quantifier is intended to range over all cluster algebras in the source's convention.

\textbf{Unresolved scope.} The displayed definition of \(\operatorname{Cl}\) applies to Krull domains. The record neither imposes that hypothesis on every algebra in the conjecture nor defines a divisor class group for the remaining algebras. Thus it does not determine a fully specified unrestricted assertion. Restricting the quantifier to Krull cluster algebras would give a precise assertion under the displayed definition, but would narrow the recorded conjecture.

The supplied dated curation labels the problem partially solved and reports that torsion-freeness is proved for acyclic cluster algebras and, more generally, for cluster algebras that are Krull domains, while the unrestricted formulation is not fully settled. These are supplied literature assessments, not independently verified mathematical or current-status claims.
\subsection{Short English statement}
The recorded conjecture asks whether every element of finite order in the divisor class group of any cluster algebra is zero. The coefficient conventions and the meaning of the class group outside the Krull-domain setting are unspecified. Supplied dated curation reports results for acyclic and Krull cluster algebras and labels the unrestricted problem partially solved; this draft requires human mathematical review.
\subsection{Sources}
\begin{itemize}
\item Supplied local UnsolvedMath registry record, permanent ID \(30002518\), source problem code \texttt{catalogue entry}, titled \emph{Torsion-Freeness of Cluster Algebra Class Groups}, in the Oberwolfach Reports Open Problems collection. The record assigns the year 2014 and names no individual proposer.
\item \emph{Representation Theory of Quivers and Finite Dimensional Algebras} (2014), Oberwolfach Reports, pp.~453--529, DOI \texttt{10.4171/owr/2014/08}. Supplied source report and literature-source label: Oberwolfach Report 08/2014. \url{https://doi.org/10.4171/owr/2014/08}. The supplied background describes these reports as documenting mathematical workshops and open-problem discussions; no exact page for this conjecture is supplied.
\item Garcia Elsener--Lampe--Smertnig, supplied literature-source attribution. \url{https://arxiv.org/abs/1712.06512}.
\item \emph{Class groups of upper cluster algebras}, supplied literature-source title. \url{https://arxiv.org/abs/2601.07520}.
\item Literature assessment checked on \texttt{2026-08-21}, exactly as supplied in the record. This is supplied dated curation derived from the OWR curation, not independently verified. It reports torsion-freeness for acyclic cluster algebras and, more generally, Krull cluster algebras, and describes the unrestricted formulation as not fully settled. The background explicitly calls for independent verification of time-sensitive status and cited claims.
\item The supplied registry labels the statement \texttt{corrected\_verified} and describes a manual semantic, typographic, and LaTeX review of its cleanup. This records the supplied editorial assessment; it does not establish mathematical correctness or resolve the scope qualifications above.
\end{itemize}
\end{document}
